\documentclass{article}
\usepackage[utf8]{inputenc}
\usepackage[T1]{fontenc}
\usepackage{ngerman}

\begin{document}

\section*{A.1. Questions Energy Systems}

\subsection*{Global Energy Resources}
\begin{enumerate}
    \item \textbf{The energy flow from mining to consumption is categorized in five stages. List these stages.} \\
    Primary Energy, Transformation Process, Secondary Energy, Final Energy Consumption, Useful Energy.
    \item \textbf{Give an example for a Primary Energy / Secondary Energy.} \\
    Primär: Rohöl, Kohle, Erdgas. Sekundär: Elektrizität.
    \item \textbf{When talking about energy sources their availability is split into Reserve and Resource. Define the meaning of these two terms Reserve and Resource.} \\
    Reserve: Mit heutiger Technologie wirtschaftlich förderbar. Ressource: Potenziell vorhanden, aber derzeit nicht wirtschaftlich abbaubar.
    \item \textbf{Name the three mostly used primary energy sources.} \\
    Rohöl, Kohle, Erdgas.
    \item \textbf{Name three renewable energy sources.} \\
    Biomasse, Wasserkraft, Windkraft (auch Solar oder Geothermie).
    \item \textbf{Name three non-renewable energy sources} \\
    Fossile Brennstoffe (Kohle, Gas, Öl), Uran.
    \item \textbf{What does the abbreviation OPEC stays for?} \\
    Organization of Petroleum Exporting Countries.
    \item \textbf{Why does the shale oil production received a boost in recent years?} \\
    Technologische Fortschritte ermöglichen die wirtschaftliche großflächige Erschließung aus dichtem Gestein.
\end{enumerate}

\subsection*{Energy Transformation}
\begin{enumerate}
    \setcounter{enumi}{8}
    \item \textbf{Name and describe the levels of a traditional hierarchical power system structure.} \\
    Produktion (Hochspannung) $\rightarrow$ Übertragung (Hochspannung) $\rightarrow$ Verteilung (Mittel-/Niederspannung) an den Endverbraucher.
    \item \textbf{Why does the hierarchical power system has several voltage levels?} \\
    Um große Energiemengen mit minimalen Leistungsverlusten über große Distanzen zu transportieren.
    \item \textbf{What is a typical voltages for LV, MV, HV voltage levels?} \\
    HV: bis 380 kV, MV: 10 bis 30 kV, LV: 400 V und darunter.
    \item \textbf{What are the abbreviations DSO and TSO staying for and give one example each?} \\
    TSO: Transmission System Operator (z.B. APG oder Vorarlberger Uebertragungsnetz). DSO: Distribution System Operator (Im Skriptum wird kein spezifisches DSO-Beispiel explizit benannt).
\end{enumerate}

\subsection*{Electrical Energy Production}
\begin{enumerate}
    \setcounter{enumi}{12}
    \item \textbf{Describe the terms base load, mid-merit load and peak load.} \\
    Grundlast (dauerhafter Bedarf), Mittellast (variabler Bedarf), Spitzenlast (kurzzeitige Leistungsspitzen).
    \item \textbf{What are typical base load power plants and what is their common characteristic?} \\
    Laufwasserkraftwerke, Nuklear- und Braunkohlekraftwerke. Charakteristik: Laufen kontinuierlich, schwer regelbar.
    \item \textbf{What is a typical mid-merit power plant?} \\
    Steinkohlekraftwerke.
    \item \textbf{What are typical peak power plants and what is their common characteristic?} \\
    Pumpspeicherkraftwerke und Gasturbinen. Charakteristik: Sehr schnelle Startzeiten.
    \item \textbf{Describe the term spinning reserve and what is it used for?} \\
    Sofort verfügbare Erzeugungskapazität durch Generatoren, die unter Volllast laufen. Dient dem Ausgleich plötzlicher Lastwechsel oder Ausfälle.
    \item \textbf{What does the Wilians line describes?} \\
    Das Verhältnis von Wärmezufuhr (GJ/h) zur abgegebenen Leistung (MW) bei fossilen Kraftwerken.
    \item \textbf{Describe the three levels of the electrical energy market.} \\
    1: Primärenergiemarkt (Produktion); 2: Kommerziell (Handel/Verkauf) sowie physikalisch (Übertragung/Verteilung); 3: Verbrauch.
    \item \textbf{Which levels of the electrical energy markets are partly regulated (natural monopole) and which are open for competition?} \\
    Übertragung und Verteilung sind reguliert. Produktion, Handel, Vertrieb und Verbrauch sind für den Wettbewerb offen.
    \item \textbf{Worldwide we see the trend of unbundling former monopolistic energy companies... What is the motivation, driving force behind this development?} \\
    Mehr Wettbewerb, Effizienzsteigerung der Industrie und Wahlfreiheit für Kunden.
    \item \textbf{What are the biggest differences between regular economic market and the electricity market?} \\
    Elektrische Energie ist in großen Mengen kaum speicherbar und der Leistungsfluss folgt physikalischen Gesetzen.
    \item \textbf{Who is responsible for the regulation in Austria?} \\
    E-Control.
    \item \textbf{Why is the tradable product Electrical Energy time based?} \\
    Wegen der fehlenden großflächigen Speicherbarkeit muss Erzeugung und Verbrauch zeitgleich erfolgen.
    \item \textbf{Why is a central coordination of the electrical infrastructure needed?} \\
    Zur Wahrung der Netzbalance (Frequenz/Spannung), zur Einhaltung starrer Sicherheitsgrenzen und weil Übertragung ein Monopol darstellt.
    \item \textbf{What are the tasks of a control area manager?} \\
    Netzüberwachung, Ausgleich von Disbalancen, Beschaffung von Systemdienstleistungen.
    \item \textbf{Name the control area managers for Austria, Germany and Switzerland.} \\
    AT: APG. DE: Amprion, 50 Hertz, TransBW, Tennet. CH: Swissgrid.
    \item \textbf{What voltage levels has the Austrian high voltage grid?} \\
    Bis zu 380 kV (inkl. 220 kV, 110 kV).
\end{enumerate}

\subsection*{Power Generation Costs}
\begin{enumerate}
    \setcounter{enumi}{28}
    \item \textbf{What are CAPEX and what are they used for?} \\
    Capital Expenditures (Investitionskosten). Sie decken den Bau und die Errichtung der Anlage ab.
    \item \textbf{Give examples for CAPEX’s.} \\
    Anlagenkosten, Komponenten, Planung, Netzanschluss.
    \item \textbf{What are OPEX and what are they used for?} \\
    Operational Expenditures (Betriebskosten). Sie fallen für den laufenden Betrieb (z.B. Brennstoffe) an.
    \item \textbf{What is the effective operating time of a power plant?} \\
    Auslastung (Utilization) multipliziert mit der Verfügbarkeit (Availability) mal 8760 Stunden.
    \item \textbf{How is the Utilization (dt. Auslastung) of a power plant defined?} \\
    Tatsächlich produzierte Energie geteilt durch die Energiemenge, die bei ständiger Laufbereitschaft unter Volllast produziert worden wäre.
    \item \textbf{How is the power generation costs $\epsilon$ defined?} \\
    Jährliche Gesamtkosten geteilt durch die jährlich produzierte elektrische Energie.
\end{enumerate}

\subsection*{Generation System Planning}
\begin{enumerate}
    \setcounter{enumi}{34}
    \item \textbf{What is the accumulated daily load curve used for?} \\
    Um die Lastsituation zu definieren und so die richtigen Kraftwerkstechnologien sowie deren minimale Betriebsstunden zu planen.
\end{enumerate}

\subsection*{Energy Trading}
\begin{enumerate}
    \setcounter{enumi}{35}
    \item \textbf{For which primary energy sources are markets existing?} \\
    Kohle, Öl, Gas, Uran.
    \item \textbf{For which primary energy sources are no markets existing?} \\
    Im Skript nicht direkt definiert (implizit Erneuerbare Energien wie Sonne/Wind).
    \item \textbf{What is traded on the Emission rights market?} \\
    Im Textausschnitt nicht explizit aufgeführt (typischerweise $CO_{2}$-Zertifikate).
    \item \textbf{What markets are available for electrical power?} \\
    Spotmärkte (Day-Ahead, Intraday) und Terminmärkte (Futures).
    \item \textbf{How can the product electrical energy be classified?} \\
    Nach Lieferperiode, Lieferprofil, Handelszeitraum und Erfüllungsart.
    \item \textbf{Product characteristics like delivery period and profile are defining the marketplace... Name two market places and the typical delivery period, delivery profile and trading period...} \\
    Spotmarkt (Periode: 1h, Profile: Base/Peak, Handel am Vortag oder Intraday). Terminmarkt/Futures (Periode: Wochen/Monate/Jahre, Profile: Base/Peak, Handel vor Lieferbeginn).
    \item \textbf{Why do have utilities have their own optimal portfolio of electrical energy products?} \\
    Um ihre Lieferverpflichtungen durch einen Mix aus Eigenproduktion und Marktgeschäften möglichst wirtschaftlich zu erfüllen.
    \item \textbf{What are the two basic ways business transactions are happening?} \\
    Kontinuierlicher Handel (Continuous trading) und Auktionen.
    \item \textbf{Which type of transaction is using an open ordering book and how does it works?} \\
    Kontinuierlicher Handel. Kauf- und Verkaufsaufträge werden gesammelt und fortlaufend abgeglichen.
    \item \textbf{Which markets are usually doing continuous trading?} \\
    Intraday-Handel.
    \item \textbf{Which type of transaction is using a closed ordering book and how does it works?} \\
    Auktionen. Aufträge werden bis zu einer Deadline gesammelt, danach wird der einheitliche Preis gebildet.
    \item \textbf{Which market is usually doing auctions?} \\
    Day-Ahead-Handel und der Handel mit Übertragungsnetzkapazitäten.
    \item \textbf{What is a Market Clearing Price?} \\
    Der Schnittpunkt der aggregierten Angebots- und Nachfragekurven; ein für alle Teilnehmer gültiger Preis.
    \item \textbf{How are business transactions being fulfilled after the product electrical energy has been delivered?} \\
    Physische Erfüllung, Cash-Settlement, Cascading, Lieferung des Futures.
    \item \textbf{What does a physical fulfillment of a business transaction means and for what electricity products is it usually used?} \\
    Tatsächliche Stromlieferung gegen Geld. Genutzt bei Spotmärkten (Day-Ahead/Intraday) und Zukunftsoptionen.
    \item \textbf{What does Cash-settlement of a business transaction means} \\
    Rein finanzieller Ausgleich der Preisdifferenz zwischen dem Future-Preis und dem Spot-Durchschnittspreis.
    \item \textbf{What does a Cascaded bushiness transaction means and for what electricity products is it usually used?} \\
    Langfristige Lieferpflichten werden in mehrere kurzfristige unterteilt. Genutzt für Quartals- oder Jahres-Futures an der Börse.
    \item \textbf{How are Options, traded on stock markets, fulfilled?} \\
    Durch die tatsächliche Lieferung des definierten Futures.
    \item \textbf{What does the abbreviation OTC trading stays for and what does it means?} \\
    Over-the-counter. Es bedeutet den direkten Handel außerhalb der Strombörse.
    \item \textbf{What is the role of a broker within an OTC business?} \\
    Sie bringen als Dienstleister Angebot und Nachfrage zusammen, ohne selbst zu handeln.
    \item \textbf{Name three power exchange trading places in Europe, what products are traded there and which areas are covered.} \\
    EXAA (Spot; AT/DE), EEX (Futures; DE/FR), EPEX Spot (Spot; DE/FR/AT/CH).
    \item \textbf{What is the task of the European Commodity Clearing (ECC), the central counter part at EPEX?} \\
    Nominierung des Stromaustauschs beim TSO und Abwicklung der Zahlungen.
    \item \textbf{Which is the most common way to buy/sell electrical energy?} \\
    Der OTC-Handel.
    \item \textbf{How does the Day Ahead Trading system works?} \\
    Aufträge werden bis 12:00 Uhr gesammelt. Angebot und Nachfrage werden aggregiert und per Auktion der Market Clearing Price ermittelt.
    \item \textbf{What is the main decision factor when a utility has to decide to produce electrical energy by them self or buying it from the market...} \\
    Der Market Clearing Price im Vergleich zu den eigenen Grenzkosten (Produktionskosten).
    \item \textbf{What does the Merit order curve shows?} \\
    Kraftwerke sortiert nach ihren marginalen Produktionskosten (von günstig zu teuer).
    \item \textbf{What is a block order?} \\
    Mehrere Aufträge werden gekoppelt und nur zusammen ausgeführt, z.B. um unwirtschaftliche Einzelstunden-Starts von Kraftwerken zu vermeiden.
    \item \textbf{What does the abbreviation Phelix-Index stays for?} \\
    Physical Electricity Index.
    \item \textbf{How does the Intraday Trading system works?} \\
    Fortlaufender Handel (offenes Orderbuch), um kurzfristige Prognoseabweichungen am aktuellen Tag auszugleichen.
    \item \textbf{What is the Futures Trading system and what is it used for?} \\
    Handel von festen Mengen zu einem Fixpreis für die Zukunft, genutzt für langfristige Preisabsicherung und Planung.
    \item \textbf{Who is in charge for trading transmission line capacities?} \\
    Die Transmission System Operators (TSO).
    \item \textbf{What is a balance group used for and who can participate in such a group?} \\
    Bündelt Erzeuger, Verbraucher und Händler zum Bilanzabgleich. Jeder Netznutzer (Einspeiser, Entnehmer, Händler) muss Mitglied sein.
    \item \textbf{Who is the central clearing and settlement agent (CSA) (dt. Bilanzgruppenkoordinator) in Austria?} \\
    Austrian Power Clearing \& Settlement (APCS).
    \item \textbf{Give three examples for balance responsible parties (dt. Bilanzgruppenverantwortliche) registered in Austria?} \\
    Vorarlberger Kraftwerke AG, ÖBB Infrastruktur, Next Kraftwerke AG.
\end{enumerate}

\subsection*{Thermal Plants}
\begin{enumerate}
    \setcounter{enumi}{69}
    \item \textbf{Name for fuel types used in thermal power plants.} \\
    Kohle, Gas, Öl, Uran (und Biomasse).
    \item \textbf{What is a cycle process (dt. Kreiprozess)?} \\
    Ein thermodynamischer Prozess, bei dem das Arbeitsmedium nach mehreren Zustandsänderungen wieder den Ausgangszustand erreicht.
    \item \textbf{Name to typically used working materials in cycle processes?} \\
    Dampf und Gas.
    \item \textbf{Describe the four stages of the Carnot cycle.} \\
    1-2 Isotherme Expansion (Wärmezufuhr), 2-3 Isentrope Expansion (Arbeitsabgabe), 3-4 Isotherme Kompression (Wärmeabfuhr), 4-1 Isentrope Kompression (Arbeitszufuhr).
    \item \textbf{Draw a typical p-v diagram of an Carnot cylcle for an ideal gas.} \\
    (Beschreibung gem. Skript:) Eine geschlossene Kurve, gebildet aus zwei Isothermen und zwei Isentropen.
    \item \textbf{Draw a typical T-s diagram of an Carnot cycle for an ideal gas.} \\
    (Beschreibung gem. Skript:) Ein Rechteck, begrenzt durch horizontale Isothermen und vertikale Isentropen.
    \item \textbf{Name the typical components of a power plant in which the Rankine cycle is realized.} \\
    Speisewasserpumpe, Dampfkessel, Überhitzer, Turbine, Kondensator.
    \item \textbf{Describe the four stages of the Rankine cycle.} \\
    1-2 Isentrope Kompression, 2-3 Isobare Wärmezufuhr/Verdampfung, 3-4 Isentrope Entspannung, 4-1 Isobare Kondensation.
    \item \textbf{Draw a typical T-s diagram of an ideal Rankine cylcle.} \\
    (Beschreibung gem. Skript:) Glockenkurve des Wassers mit isentroper Pumparbeit, isobarer Verdampfung, isentroper Expansion und isobarer Kondensation.
    \item \textbf{What is the difference between an ideal and an real Rankine process.} \\
    Der reale Prozess hat Irreversibilitäten: Druckabfälle und nicht-isentrope (verlustbehaftete) Kompression/Expansion in Pumpe und Turbine.
    \item \textbf{What is the dryness fraction x (dt. Dampfgehalt) used for?} \\
    Er beschreibt im Nassdampfgebiet das Massenverhältnis von Sattdampf zur gesamten Flüssig-Dampf-Mischung.
    \item \textbf{What are the two extreme points of the dryness fraction x describing? (x=0, x=1)} \\
    x=0: reines flüssiges Wasser (Siedelinie); x=1: reiner Sattdampf (Taulinie).
    \item \textbf{What are the three aggregation states of the working medium steam?} \\
    Flüssigkeit, Nassdampf (2-Phasen-Gemisch), überhitzter Dampf.
    \item \textbf{Name two application areas of gas turbines.} \\
    Flugzeugantriebe und Stromerzeugung.
    \item \textbf{What is a typical start-up time of a gas turbine?} \\
    Ca. 2 Minuten.
    \item \textbf{Name the typical components of a gas turbine where the Bryton cycle is realized.} \\
    Lufteinlass, Kompressor, Brennkammer, Turbine, Auspuff.
    \item \textbf{Describe the four stages of the Brayton cycle.} \\
    1-2 Isentrope Kompression, 2-3 Isobare Erwärmung, 3-4 Isentrope Entspannung, 4-1 Isobare Wärmeabfuhr.
    \item \textbf{Draw a typical T-s diagram of an ideal Brayton cylcle.} \\
    (Beschreibung gem. Skript:) Eine geschlossene Kurve zwischen zwei vertikalen Isentropen und zwei Isobaren.
    \item \textbf{Name the components of a coal fired power plant.} \\
    Feuerung, Kohlemühle, Frischluftgebläse, Kessel, Turbine, Kondensator, DeNOx-Anlage, Elektrofilter, Saugzuggebläse, Entschwefelung.
    \item \textbf{Describe a typical steps of a fixed bed or grate firing process.} \\
    Brennstoffaufgabe, Trocknung, Zündung, Verbrennung, Ascheabfuhr.
    \item \textbf{What are the advantages / disadvantages of a grate firing system with a feed stoker?} \\
    Vorteil: Einfach, gute Lastwechsel, weniger NOx. Nachteil: Unverbrannte Verluste, begrenzte Luftvorwärmung, ungeeignet für Feinkohle.
    \item \textbf{What is a Bubbling fluidized bed?} \\
    Stationäre Wirbelschicht mit niedriger Gasgeschwindigkeit zur Fluidisierung.
    \item \textbf{What is a Circulating fluidized bed (CFB)?} \\
    Zirkulierende Wirbelschicht mit hoher Gasgeschwindigkeit, Partikel werden extern rezirkuliert.
    \item \textbf{What are the advantages / disadvantages of dust firing?} \\
    Vorteil: Schneller Start, schnelle Lastwechsel. Nachteil: Hoher Mahl-Energiebedarf, mehr Emissionen, aufwendige Wartung.
    \item \textbf{Name three ways a coal burner can be installed in a boiler.} \\
    Einseitig, gegenüberliegend, tangential (in den Ecken).
    \item \textbf{What is the difference between primary and secondary air used by an coal burner?} \\
    Primärluft transportiert die Kohle, Sekundärluft wird zur eigentlichen Verbrennung eingeblasen.
    \item \textbf{Name two methods on how the initial ignition of a coal burner is done.} \\
    Ölbrenner oder Lichtbogen zwischen Elektroden.
    \item \textbf{What are the two basic different designs of a steam generator?} \\
    Naturumlauf und Zwangsumlauf.
    \item \textbf{What does the parameter circulation ratio describes within a steam generator?} \\
    Das Verhältnis zwischen der umlaufenden Wassermenge und der produzierten Dampfmenge.
    \item \textbf{What are condenser used for in a thermal power plant?} \\
    Um den thermodynamischen Kreisprozess zu schließen, indem Abluftdampf der Turbine wieder in Speisewasser kondensiert wird.
    \item \textbf{How can condenser be classified?} \\
    Mischkondensatoren und Oberflächenkondensatoren.
    \item \textbf{Explain the working principle of Electrostatic precipitator (ESP).} \\
    Rauchgas strömt durch ein DC-Hochspannungsfeld. Partikel werden ionisiert, von Elektroden angezogen und entfernt.
    \item \textbf{Describe how a wet scrubber is used for flue gas desulphurization.} \\
    Abgas strömt aufwärts, während Kalkwasser von oben herabrieselt und durch Reaktion das $SO_{2}$ aus dem Gas wäscht.
    \item \textbf{Name two methods for $NO_{x}$ emission control.} \\
    Selective catalytic reduction (SCR) und Selective non-catalytic reduction (SNCR).
    \item \textbf{Name the components of a steam turbine.} \\
    Leitrad (Düsen), Laufrad (Rotor), Gehäuse.
    \item \textbf{What is a turbine stage?} \\
    Die Kombination aus Leitrad und Laufrad.
    \item \textbf{Describe the Reaction Principle of a steam turbine.} \\
    Dampf entspannt sich im Leit- und Laufrad (Druckabfall in beiden), Reaktionskräfte treiben den Rotor an.
    \item \textbf{Describe the Impulse Principle of a steam turbine.} \\
    Dampf entspannt vollständig im Leitrad zu einem Schnellstrahl. Im Laufrad wird kinetische Energie ohne Druckabfall übertragen.
    \item \textbf{What is the working principle of combined-cycle gas-turbine (CCGT) plants?} \\
    Die heißen Abgase einer Gasturbine werden genutzt, um in einem Kessel Dampf für eine zweite Dampfturbine zu erzeugen.
    \item \textbf{What are the feature of a CCGT?} \\
    Sehr hoher Wirkungsgrad ($\sim60\%$), schnelle Starts, flexibler Einsatz (auch mit Leichtöl).
    \item \textbf{Where does the heat in a nuclear power plant comes from?} \\
    Aus der Kernspaltung (Fission) von Uranatomen.
    \item \textbf{What does the term moderator means in a nuclear power plant?} \\
    Ein Medium, das Neutronen abbremst, um eine effektive Kernspaltung aufrechtzuerhalten.
    \item \textbf{What do the terms light water and heavy water means in a nuclear power plant?} \\
    Light water = konventionelles $H_{2}O$. Heavy water = Deuteriumoxid ($D_{2}O$).
    \item \textbf{Describe boiling water reactor (BWR) and Pressurized Water Reactor (PWR).} \\
    PWR: Wasser steht unter Hochdruck (siedet nicht) und gibt Wärme an einen Sekundärkreis ab. BWR: Wasser siedet direkt im Reaktor.
    \item \textbf{Name the two main tasks of a thermal power plant controller.} \\
    Kesselregelung (Boiler control) und Turbinenregelung (Turbine control).
    \item \textbf{Describe the control strategy Boiler following.} \\
    Turbinenventil regelt Last direkt. Der entstehende Druckabfall veranlasst den Kessel, die Feuerung nachzuregeln.
    \item \textbf{Describe the control strategy Turbine following.} \\
    Kessel regelt Feuerung primär; die Turbine passt ihr Ventil an den neuen Dampfdruck an.
\end{enumerate}

\subsection*{Renewable Plants}
\begin{enumerate}
    \setcounter{enumi}{116}
    \item \textbf{How can hydro power plants be categorized?} \\
    Laufwasserkraftwerke (River power plants), Speicherkraftwerke, Ozeankraftwerke.
    \item \textbf{What is the advantages of using a diversion canal for river power plants?} \\
    Im bereitgestellten Text wird kein spezifischer Vorteil für Ausleitungskanäle explizit genannt.
    \item \textbf{Why is the power house placed on the outer river bank on a bending river?} \\
    Im bereitgestellten Text nicht explizit erklärt.
    \item \textbf{What is the commonly used turbine for river power plants, and why?} \\
    Kaplan-Turbinen, da sie sehr gut für niedrige Fallhöhen (bis 25m) geeignet sind.
    \item \textbf{What load type (base, medium, peak) is usually covered by river power plants?} \\
    Grundlast (Base load).
    \item \textbf{What is a commonly used turbine used for storage power plants, and why?} \\
    Francis-Turbinen (hohe Wassermengen, 15-500m) und Pelton-Turbinen (geringere Mengen, bis 2000m Fallhöhe).
    \item \textbf{What are cascaded river power plants and what possible operation modes are used to operate them?} \\
    Eine Reihe hintereinander liegender Flusskraftwerke (Staukette). Modi: Kippbetrieb, Schwellbetrieb.
    \item \textbf{What is the difference between the tipping operation and swell operation of cascaded river power plants?} \\
    Kippbetrieb: Alle Kraftwerke öffnen gleichzeitig. Schwellbetrieb: Kraftwerke aktivieren die Produktion sequenziell wie eine Welle.
    \item \textbf{Describe a typical cyclic operation (dt. Wälzbetrieb) of an pump storage power plant.} \\
    Wasser fällt zur Spitzenlastzeit abwärts und generiert Strom; in Schwachlastzeiten (nachts) wird es hochgepumpt.
    \item \textbf{Name at least 3 objectives of a pump storage power plant.} \\
    Spitzenlastabdeckung, Veredelung von Grundlast, Frequenz-/Spannungsregelung (Netzstützung), Schwarzstartfähigkeit.
    \item \textbf{How are combined pump storage power plants defined?} \\
    Die Pumpe verarbeitet weniger als die Hälfte des Wassers, das durch die Turbine fließt (natürlicher Zufluss überwiegt).
    \item \textbf{What is a tandem set (dt. Tandemsatz) in a pump storage power plant?} \\
    Ein System mit komplett getrennter Pumpe und Turbine auf einer Welle.
    \item \textbf{Describe what the operation mode hydraulic short circuit means and what it is used for.} \\
    Im Text nicht explizit definiert.
    \item \textbf{What is a typical value for a pump storage efficiency?} \\
    Das Skript beschreibt sie als „deutlich besser als thermische Kraftwerke“ (der genaue %-Wert fehlt im vorliegenden Text).
    \item \textbf{Name three types of oceanic power plants.} \\
    Im bereitgestellten Text werden keine 3 Typen explizit benannt.
    \item \textbf{Name three characteristics of water.} \\
    Dichteanomalie (+4°C bis -4°C), siedet im Unterdruck bei sehr niedrigen Temperaturen, universelles Lösungsmittel.
    \item \textbf{What are the two parameters the power output of a hydro power plant depends on the most?} \\
    Wassermasse (Durchfluss) und Fallhöhe.
    \item \textbf{Describe the phenomenon cavitation erosion.} \\
    Wasserdampfblasen entstehen durch Unterdruck. Bei erneutem Druckanstieg implodieren diese und schädigen Material.
    \item \textbf{Name the three most important hydro turbines.} \\
    Francis-Turbine, Pelton-Turbine, Kaplan-Turbine.
    \item \textbf{Describe how a Francis turbine works.} \\
    Reaktionsturbine: Wasser strömt radial ein, wird um 90° umgelenkt und verlässt den Rotor axial. Druckabfall treibt Rotor.
    \item \textbf{Describe how a Pelton turbine works.} \\
    Gleichdruckturbine: Hochdruck-Wasserstrahlen treffen auf Becher, teilen sich und ändern die Richtung, was einen Impuls erzeugt.
    \item \textbf{Describe how a Kaplan turbine works.} \\
    Propellerähnlicher Rotor mit verstellbaren Schaufeln zur Anpassung an Durchfluss und Höhe (Reaktionsprinzip).
    \item \textbf{Describe how photovoltaic conversion works.} \\
    Photonen erzeugen Elektron-Loch-Paare; deren Trennung generiert eine elektrische Spannung.
    \item \textbf{Name the components of a wind turbine?} \\
    Rotor (Blätter/Nabe), Antriebsstrang (Getriebe/Generator/Bremse), Nachführsystem (Yaw), Turm/Fundament, Leistungselektronik.
    \item \textbf{What makes the blades of a wind turbine move? Describe the physical principle.} \\
    Aerodynamische Auftriebskräfte durch Druckunterschiede aufgrund unterschiedlicher Strömungsgeschwindigkeiten über das Blattprofil.
    \item \textbf{Name three typically used generators in wind turbines?} \\
    Direktgekoppelte Asynchrongeneratoren, doppeltgespeiste Asynchrongeneratoren (DFIG), Synchrongeneratoren.
    \item \textbf{Describe the working principle of an doubly fed induction generator (dt. DoppelGespeisten Asynchongenerator).} \\
    Stator ist starr am Netz; der Rotor wird über einen Umrichter manipuliert, wodurch variable Drehzahlen möglich sind.
    \item \textbf{What are the so called Grid Codes?} \\
    Vorgaben der Übertragungsnetzbetreiber (TSOs), wie sich Anlagen am Netz verhalten müssen (z.B. Fault-Ride-Through, Frequenzstützung).
    \item \textbf{Name two methods which are used to limit power capture if the wind are getting to strong.} \\
    Pitch-Regelung, Stall-Regelung.
    \item \textbf{A traditional way of turbine control is the so called Danish Concept. What does it consists of?} \\
    Passive Stall-Regelung und starrer Direktanschluss ans Netz via Getriebe und Asynchrongenerator.
    \item \textbf{What is the advantage of the variable speed control concept?} \\
    Bessere Regelbarkeit, weniger mechanische Belastung, verbesserte Netzqualität und höhere Energieausbeute.
    \item \textbf{Name the two ways biomass can be used to produce electricity.} \\
    Im Text (S. 130) ist nur „Biofuels derived from vegetable matter...“ beschrieben (Text bricht ab). Zuvor wird feste Biomasse (Holz) genannt.
\end{enumerate}

\subsection*{Power System Control}
\begin{enumerate}
    \setcounter{enumi}{148}
    \item \textbf{What are the tasks of a control area manager (dt. Regelzonenführer)?} \\
    Überwachung der Netzbalance (Frequenz und Spannung), Zukauf und Abruf von Systemdienstleistungen (Reserveenergie).
    \item \textbf{Describe the working principle of an synchronous generator?} \\
    Ein mit Gleichstrom erregter Rotor erzeugt ein rotierendes Magnetfeld, welches im direkt ans Netz gekoppelten Stator Wechselspannung induziert.
    \item \textbf{How is grid frequency related to real power supply?} \\
    Frequenz fällt bei Leistungsunterdeckung (Generatoren werden gebremst) und steigt bei Leistungsüberschuss.
    \item \textbf{How is grid voltage related to reactive power supply?} \\
    Die Spannung wird durch die Bereitstellung oder Aufnahme von Blindleistung gestützt.
    \item \textbf{What are the three-stage procedure consists of, which is used to keep the electrical power balanced within the Continental European Grid?} \\
    Primärregelung, Sekundärregelung, Tertiärregelung.
    \item \textbf{What is the time period for a single incident in the primary control stage?} \\
    0 bis 30 Sekunden.
    \item \textbf{Who participates in the primary control?} \\
    Alle Übertragungsnetzbetreiber (TSOs) im Verbundsystem (ENTSO-E).
    \item \textbf{What is the time period for a single incident in the secondary control stage?} \\
    30 Sekunden bis 15 Minuten.
    \item \textbf{Who participates in the secondary control?} \\
    Kraftwerke im Regelgebiet, die in Betrieb sind, aber nicht unter Volllast laufen.
    \item \textbf{What is the time period for a single incident in the tertiary control stage?} \\
    Im vorliegenden Textausschnitt nicht spezifiziert (typisch: > 15 Minuten).
    \item \textbf{Who participates in the tertiary control?} \\
    Im vorliegenden Textausschnitt nicht spezifiziert.
    \item \textbf{What does the minutes reserve means?} \\
    Im vorliegenden Textausschnitt nicht definiert.
\end{enumerate}

\end{document}